\section{Introduction}
Most Tamil speakers type on phones using the Latin alphabet.
Search, autocomplete and display all need the text converted back into Tamil script, which is a transliteration task: sounds are mapped to letters, not meanings to words.
The Latin alphabet cannot express several Tamil distinctions, including vowel length, doubled consonants and the several n, l and r sounds, and writers spell the same word in different ways.
Informal typing adds more variation: dropped vowels (\emph{kankl} for \emph{kanakkil}), collapsed long vowels, collapsed doubled consonants, swapped digraphs and stretched letters.

Neural transliteration is well studied on clean data.
The Dakshina dataset provides romanizations of Tamil words written by native speakers on request~\cite{roark2020dakshina}, and strong systems have been trained on larger corpora~\cite{madhani2023aksharantar}.
However, the Tamil text in Dakshina is encyclopaedic and was typed carefully, so it does not contain the shorthand found in chat.
We therefore ask how much character-level models lose under informal typing, which training strategy recovers it, and whether the choice of architecture matters.

The literature also disagrees on architecture.
On Dakshina Tamil, an LSTM beat a Transformer in one study~\cite{roark2020dakshina}, whereas a Transformer was better in another~\cite{kirov2024context}, and batch size can reverse the ranking~\cite{wu2021applying}.
None of these comparisons tested noisy input.

\textbf{Contributions.}
(1)~A paired, seeded robustness benchmark for Tamil with seven noise operators and four severities.
(2)~A study of four augmentation strategies and a more-data control, evaluated on held-out noise types and on real typing.
(3)~A parameter-matched, equal-tuning comparison of BiLSTM-attention and Transformer over five seeds with paired significance tests.
(4)~A reproducibility package: code, configurations, per-run logs and data checksums.
